\documentclass[10pt,a4paper]{amsart}
\usepackage[margin=19mm]{geometry}
\usepackage{amsmath,amssymb,mathtools}
\usepackage[scr=rsfs,cal=euler]{mathalfa}
\usepackage{xfrac}
\usepackage[unicode,hidelinks,bookmarks=false]{hyperref}
\newcommand{\ie}{i.e.\ }
\newcommand{\cf}{cf.\ }
\newcommand{\resp}{respectively\ }
\newcommand{\Subsection}{\subsection}
\providecommand{\nobreakdash}{\nobreak\hbox{-}}
\setlength{\emergencystretch}{2em}
\multlinegap0pt
\newtheorem{lemma}{Lemma}
\newcommand{\Div}{{\rm div}}
\newcommand{\Th}{\mathcal{T}_h}
\newcommand{\av}[1]{\{\kern-1.2mm \{{#1}\}\kern-1.2mm\}}
\newcommand{\avg }[1]{ \av{ #1}}
\newcommand{\cinv}{C_{\nabla}}
\newcommand{\cinvT}{C_{\partial}}
\newcommand{\dee}{{\rm d}}
\newcommand{\dt}{\,{\rm dt}}
\newcommand{\gif}{\Gamma_-}
\newcommand{\gint}{\Gamma_{\rm int}}
\newcommand{\inner}[3]{({#1},{#2})_{{#3}}}
\newcommand{\jump}[1]{[\kern-0.6mm [{#1}]\kern-0.6mm]}
\newcommand{\jumpu}[1]{ \lfloor #1 \rfloor }
\newcommand{\norm}[2]{\|{#1}\|_{{#2}}}
\newcommand{\normal}{{n}}
\newcommand{\su}{\sum_{T \in \Th}}
\newcommand{\tf}{{t_{\rm f}}}
\newcommand{\utr}{W}
\newcommand{\xn}{x_{n_2}^{T}}
\title{Asymptotic numerical hypocoercivity \\ of the space-time discontinuous Galerkin method for {the} Kolmogorov equation}
\author{Zhaonan Dong, Emmanuil H. Georgoulis, Philip J. Herbert}
\date{Source: arXiv:2505.19222v3 (CC BY 4.0)}
\begin{document}
\maketitle

\noindent\emph{Source and licence.} Asymptotic numerical hypocoercivity \\ of the space-time discontinuous Galerkin method for {the} Kolmogorov equation. Version arXiv:2505.19222v3 (CC BY 4.0), distributed by arXiv under the Creative Commons Attribution 4.0 International licence, http://creativecommons.org/licenses/by/4.0/, as recorded in the arXiv licence metadata for this exact version and shown on the abstract page https://arxiv.org/abs/2505.19222v3. Full licence notice: \texttt{licenses/arxiv-2505.19222-CC-BY-4.0.txt}.\par\smallskip

\noindent\textit{Editorial scope.} Counted unit: the article's lemma lem:semi:hypo (source0.tex lines 708--732). The article's complete original proof of it is delivered with it (source0.tex lines 733--911, one \texttt{proof} environment of 179 lines). No mathematics is written, repaired or re-derived: the statement and the proof are the article's own lines, in the article's own notation, and no retained line is substituted or re-flowed.  Retained uncounted as the source-local context the proof needs: lines 297--301; lines 303--312; lines 403--407; lines 432--439; lines 524--527; lines 569--576; lines 644--653; lines 1899--1907.  Nothing was omitted from any retained span, so the package carries no omission record.  No retained line contains a citation, so the wrapper carries no bibliography and no bibliography entry.  No retained line contains a figure environment, an image-inclusion command or drawing code, so no image is delivered and the package declares no asset dependency.  Editorial formatting only, in addition: an AMS \texttt{amsart} preamble that restates the article's own declarations and macros used by the retained lines, namely the environments \texttt{lemma} on separate counters, together with the standard packages those lines need.  The retained environments print in this document as Lemma 1 in order of appearance, whereas the article prints the same items with the numbering of its own full document.  The complete native source of the article as distributed in the arXiv e-print for this exact version is bundled unchanged as \texttt{sources/178739/source0.tex}.
\begin{equation}\label{eq:A_def}
 A : =      \begin{pmatrix}
     \alpha & \beta \\ \beta & \gamma
    \end{pmatrix}
\end{equation}

\begin{equation}\label{eq:alpha_beta_gamma_def}
    \alpha =
    \frac{1}{8\delta},
    \quad
    \beta =
    \frac{1}{48 \delta^2},
    \quad
    \gamma =
    \frac{10}{48^2 \delta ^3}
\end{equation}

\begin{equation}\label{trace_H1_inv}
\|V\|_{\partial T} \leq \frac{\cinvT p}{ \sqrt{h_T}}  \|V\|_{ T},
 \quad\text{and}\quad
\|\nabla V\|_{T} \leq  \frac{\cinv p^2}{h_T}  \|V\|_{ T},
\end{equation}

\begin{equation}\label{def: seim-discrete bilinear form}
    B_{\rm s}(U,V)
    :=
    \int_0^{\tf}  \left(\inner{U_t}{V}{} +  a_{\rm dG}(U,V) \right) \dt
    =
    \int_0^\tf
    \inner{f}{V}{}=:   \ell_{\rm s}(V)
    \end{equation}

\begin{equation}\label{def:delta_def}\begin{split} \delta_T:=&\ \max\bigg\{1,\  \frac{\cinvT^2p^4}{h_T^2}\max\Big\{4+\frac{\cinv^2}{8
    	\cinvT^2},\frac{(\xn)^2h_T}{p^2},\sqrt{\mathcal{R}_T}\Big\}\bigg\}, \quad \text{for all $T\in \Th$}.
    \end{split}
\end{equation}

\begin{equation}\label{eq:standard_flux_estimate}
\begin{aligned}
 \inner{\avg{\normal_1 U_x}}{\jump{V_1}}{\gint }
 &\leq  \Big(\sum_{F\in \gint } \phi^{-1}(\cinvT n_1^Fp)^2 \sum_{T\in\omega_F}   h_T^{-1} \norm{  U_x}{ T}^2\Big)^{\frac{1}{2}} \norm{\sqrt{\phi} \jump{V_1}}{\gint} \\
 &\leq  \Big(\su  \norm{  U_x}{ T}^2 \Big)^{\frac{1}{2}} \norm{\sqrt{\phi} \jump{V_1}}{\gint} \\
&\leq \frac{1}{2\varepsilon}  \su \norm{ U_x}{ T}^2 +  \frac{\varepsilon}{2}\norm{\sqrt{\phi} \jump{V_1}}{\gint}^2,
\end{aligned}
\end{equation}

\begin{equation}\begin{split}\label{eq:double_inv_flux_estimate}
		\inner{\avg{\normal_1 U_x}}{\jump{\tau\utr}}{\gint}
		\ge&
		-
		\frac{1}{32}\su \norm{U_x}{T}^2 -8 \norm{\sqrt{\phi}\jump{\tau\utr}}{\gint}^2\\
		\ge &	-
		\frac{1}{32}\su \norm{U_x}{T}^2-16\sum_{F\in \gint }\! (\cinvT n_1^Fp)^2 \!\sum_{T\in\omega_F}  \! h_T^{-1} \norm{ \sqrt{\phi_T} \tau \utr}{ T}^2\\
			\ge &	-
		\frac{1}{32}\su \norm{U_x}{T}^2-16\su \phi_T^2\norm{  \tau \utr}{ T}^2,
\end{split}\end{equation}

 \begin{lemma}\label{eigen_est}
 	Consider a real $2\times 2$ symmetric and positive-semidefinite matrix $\scriptsize\mathcal{A}=\begin{pmatrix}
 		a & b \\
 		b & c
 	\end{pmatrix},
 $ whose entries satisfy
$a,c\ge0$. Then, we have $
\frac{ac-b^2}{a+|b|}\le |\mathcal{A}| \le a+|b|\le 2a$, with $|\cdot|$ the spectral matrix norm.
 \end{lemma}

\begin{lemma}\label{lem:semi:hypo}
    Let $U \in H^1(I; {\mathcal{V}}_h)$, $A$ given by \eqref{eq:A_def} with $\alpha$, $\beta$, and $\gamma$ as in \eqref{eq:alpha_beta_gamma_def}, and $\delta$ as in \eqref{def:delta_def}. Then, a.e.~in $t$, we have the bound
    \begin{equation*}\begin{split}
        \inner{U_t}{V_3}{}  +    a_{\rm dG}(U,V_3)
        \geq&\,
        \frac{1}{2}  \sum_{T \in \Th}\! \Big(
       \frac{\dee}{\dt} \norm{\sqrt{A} \nabla U}{T}^2
        +
      \norm{\sqrt{A}\nabla U_x}{T}^2
        +
      \norm{ \sqrt{x\normal_2 A}\nabla U }{\partial_+ T}^2
        \\
        &
        \hspace{0cm}
        -
        \frac{1}{8} \norm{ \sqrt{\tau}(U_t + xU_y) }{T}^2
        +
        \frac{\alpha^2}{3}\norm{U_y}{T}^2
        -
        \frac{5}{4} \norm{U_x}{T}^2
        \Big)
        -
        \frac{1}{16}  \|U\|_{\rm dG}^2.
    \end{split}\end{equation*}
\end{lemma}

\begin{proof}


	Testing \eqref{def: seim-discrete bilinear form} with $V_3 = -\Div_h(A \nabla_h U)$, integrations by parts give
	\begin{equation*}\begin{split}
		&	  \inner{U_t}{V_3}{}  +    a_{\rm dG}(U,V_3)
			  \\
			    =&
			\sum_{T \in \Th} \bigg(
			\frac{1}{2} \frac{\dee}{\dt} \norm{\sqrt{A} \nabla U}{T}^2
			+
			\norm{\sqrt{A}\nabla U_x}{T}^2	+
			\inner{\nabla (x U_y)}{(A\nabla U)}{T}
			\\
			&\qquad\quad
			-
			\inner{U_t+ x U_y}{\normal \cdot (A\nabla U)}{\partial T}
			-
			\inner{U_x}{\normal \cdot (A\nabla U_x)}{\partial T}		\\
			&\qquad\quad
			+
			\inner{x\normal_2 \jumpu{U}}{ V_3^+ }{\partial_- T \cap \gint}
			+
			\inner{x\normal_2 U^+}{V_3^+}{\partial_- T \cap \gif}
			\bigg)
			\\
			&-
			\inner{\avg{\normal_1 U_x}}{\jump{ V_3}}{\gint }
			-
			\inner{\avg{\normal_1 \Div(A\nabla U_x) }}{\jump{U}}{\gint }
			+
			\inner{\sigma \jump{U}}{\jump{V_3}}{\gint }.
	\end{split}\end{equation*}
We estimate each term separately.

Focusing on the \emph{special} term $	\inner{\nabla(x U_y)}{A \nabla U}{T}$, we have, respectively,
	\begin{equation*}
		\begin{aligned}
	&	\inner{\nabla(x U_y)}{A \nabla U}{T}\\
		=&\
		 \alpha \inner{U_x}{U_y}{T}
		 +
		 \beta \norm{U_y}{T}^2
		 +
		 \inner{(x\sqrt{A} \nabla U)_y}{ \sqrt{A} \nabla U}{T}\\
		 =&\
		\alpha \inner{U_x}{U_y}{T}
		+
		\beta \norm{U_y}{T}^2
		+
		\frac{1}{2} \norm{ \sqrt{x\normal_2 A}\nabla U }{\partial_+ T}^2
		-
		\frac{1}{2} \norm{ \sqrt{|x\normal_2| A}\nabla U }{\partial_- T}^2\\
		\ge & \
	 \Big(\beta-\frac{2\alpha^2}{3}\Big)\norm{U_y}{T}^2	-	\frac{3}{8} \norm{U_x}{T} ^2+	\frac{1}{2} \norm{ \sqrt{x\normal_2 A}\nabla U }{\partial_+ T}^2
	 - \frac{(\xn\cinvT p)^2}{2h_T} \norm{\sqrt{A}\nabla U }{T}^2,
		\end{aligned}
	\end{equation*}
using the trace inverse inequality \eqref{trace_H1_inv}, with the understanding that $\alpha,\beta,\gamma$ are evaluated on $T$.


Now, from \eqref{trace_H1_inv}, along with elementary manipulation, we have
	\begin{equation*}\begin{split}
		-	   \inner{U_t+ x U_y}{\normal \cdot (A\nabla U)}{\partial T}
			\geq&\
		-	\frac{\cinvT ^2	p^2}{\sqrt{\tau_T}h_T} \norm{ \sqrt{\tau} (U_t+ x U_y)}{T}
	 \norm{A \nabla U}{ T}
			\\
			\geq &\
		-			\frac{1}{16}	\norm{ \sqrt{\tau} (U_t+ x U_y)}{T}^2 -  \frac{4\cinvT^4p^4}{\tau h_T^2}
			\norm{A \nabla U}{T}^2.
	\end{split}\end{equation*}
Working in a similar fashion, we also have
	\begin{equation*}
		-	\inner{U_x}{\normal \cdot (A\nabla U_x)}{\partial T}
			\geq
		-	\frac{1}{32}     \norm{U_x}{T}^2 -  \frac{ 8\cinvT^4p^4}{h_T^2}\norm{A\nabla U_x}{T}^2.
\end{equation*}
Next, working as in \eqref{eq:standard_flux_estimate} (for $\varepsilon=16$), and continuing as in \eqref{eq:double_inv_flux_estimate}, we have
	\begin{equation*}\begin{split}
			   \inner{\avg{\normal_1 U_x}}{\jump{V_3 }}{\gint}
			   \geq&\ - \frac{1}{32}  \su \norm{ U_x}{ T}^2 -  8\norm{\sqrt{\phi} \jump{V_3}}{\gint}^2\\
	   \geq&\ - \frac{1}{32}  \su \norm{ U_x}{ T}^2 -  16\su\phi_T^2\norm{V_3}{T}^2\\
	   \geq&\ - \frac{1}{32}  \su \norm{ U_x}{ T}^2 -  \su\frac{32\cinv^2\phi_T^2 p^4}{h_T^2}\norm{A\nabla U}{T}^2,
	\end{split}\end{equation*}
with an additional use of the inverse inequality \eqref{trace_H1_inv} in the last step.


Further, working as in \eqref{eq:standard_flux_estimate} (for $\varepsilon=4$) and recalling that $\sigma=64\phi$, we have
	\begin{equation*}\begin{split}
			 \inner{\avg{\normal_1 \Div(A\nabla U_x) }}{\jump{U}}{\gint}
			  \geq&\ -\frac{1}{8}  \su \norm{ \Div(A\nabla U_x) }{ T}^2 -  \frac{1}{32}\norm{\sqrt{\sigma} \jump{U}}{\gint}^2\\
			    \geq&\ -  \su \frac{\cinv^2p^4}{4h_T^2}\norm{ A\nabla U_x }{ T}^2 -  \frac{1}{32}\norm{\sqrt{\sigma} \jump{U}}{\gint}^2,
	\end{split}\end{equation*}
	and, similarly,
	\begin{equation*}\begin{split}
		-	   \inner{\sigma \jump{U}}{\jump{ V_3 }}{\gint }
			\geq& -\frac{1}{4}
			\su  \sigma_T^2\norm{ V_3}{T}^2
			-
			\frac{1}{32}   \norm{\sqrt{\sigma}\jump{U}}{\gint }^2\\
				\geq& -
			\su  \frac{\cinv^2\sigma_T^2p^4}{2h_T^2}\norm{A\nabla U}{T}^2
			-
			\frac{1}{32}   \norm{\sqrt{\sigma}\jump{U}}{\gint }^2.
	\end{split}\end{equation*}
	Finally, we also have
	\begin{equation*}\begin{split}
			&\su \Big( \inner{x\normal_2 \jumpu{U}}{ \Div(A\nabla U^+) }{\partial_- T \cap \gint}
			+
			\inner{x\normal_2 U^+}{ \Div(A\nabla U^+) }{\partial_- T \cap \gif}\Big)
			\\
			&\geq  -\sqrt{2}\norm{U}{\rm uw}
	\su	\frac{	\cinvT\cinv	\xn p^3}{h_T^{\frac{3}{2}}}
			\norm{ A \nabla U }{T}\\
			&
			\geq
			-\frac{1}{32}\norm{U}{\rm uw}^2
			-
			\su\frac{(4\cinvT\cinv	\xn p^3)^2}{h_T^3}
			\norm{ A\nabla U }{T}^2.
	\end{split}\end{equation*}
	For convenience, we denote by $|A|$ the spectral norm of $A$, i.e., its largest eigenvalue, and we note the elementary estimate $|A| \leq 2 \alpha=(4\delta)^{-1}$. The proof is presented in Lemma \ref{eigen_est}. Also, observing that $\max_{i=1,2}h_i^{-1}\le C_{\rho} h_T^{-1}$, and that $\sigma_T=64\phi_T$, we have
	\[
	\frac{\cinv^2\sigma_T^2p^4}{2h_T^2}+  \frac{32\cinv^2\phi_T^2 p^4}{h_T^2}\le \frac{2080N_\partial (n_1^FC_\rho\cinvT)^2p^8}{h_T^4},
	\]
	so that, the definition of $\delta$, \eqref{def:delta_def}, implies
	\[
	\Big( \Big(	   \frac{4\cinvT^4p^4}{\tau h_T^2}+ \frac{2080N_\partial (n_1^FC_\rho\cinvT)^2p^8}{h_T^4}+ \frac{(4\cinvT\cinv	\xn p^3)^2}{h_T^3} \Big)|A|_T| + \frac{(\xn\cinvT p)^2}{2h_T}\Big)\le \delta_T,
	\]
	and also
	$
	\big(1 - \big( 8\cinvT^4+   \frac{\cinv^2}{4} \big) \frac{| A|_T |p^4}{h_T^2}\big) \ge \frac{1}{2}$.


	Combining the above bounds, we arrive at
\begin{equation}\label{eq:inder_Vthree}
	\begin{aligned}
			&     \inner{U_t}{V_3)}{}  +    a_{\rm dG}(U,V_3) \\
		\geq
	&		\sum_{T \in \Th} \bigg(
		\frac{1}{2} \frac{\dee}{\dt} \norm{\sqrt{A} \nabla U}{T}^2
		+
	\frac{1}{2}  \norm{\sqrt{A}\nabla U_x}{T}^2+	\frac{1}{2} \norm{ \sqrt{x\normal_2 A}\nabla U }{\partial_+ T}^2- \delta_T\norm{\sqrt{A}\nabla U}{T}^2
		\\
	&\qquad	-			\frac{1}{16}	\norm{ \sqrt{\tau} (U_t+ x U_y)}{T}^2
		+ \Big(\beta-\frac{2\alpha^2}{3}\Big)\norm{U_y}{T}^2	-	\frac{3}{8} \norm{U_x}{T} ^2\bigg)	-\frac{1}{16}\norm{U}{\rm dG}^2.
	\end{aligned}
\end{equation}
Also, since
\[
\|\sqrt{A}\nabla U\|_T^2= \alpha \|U_x\|_T^2+ 2\beta (U_x, U_y)_T + \gamma \|U_y\|_T^2\le  2\alpha \|U_x\|_T^2+\Big( \gamma + \frac{\beta^2}{\alpha }\Big) \|U_y\|_T^2,
\]
 {using the above bound, we infer
\[
	\begin{aligned}
			&     \inner{U_t}{V_3)}{}  +    a_{\rm dG}(U,V_3) \\
		\geq
	&		\sum_{T \in \Th} \bigg(
		\frac{1}{2} \frac{\dee}{\dt} \norm{\sqrt{A} \nabla U}{T}^2
		+
	\frac{1}{2}  \norm{\sqrt{A}\nabla U_x}{T}^2+	\frac{1}{2} \norm{ \sqrt{x\normal_2 A}\nabla U }{\partial_+ T}^2- \delta_T\norm{\sqrt{A}\nabla U}{T}^2
		\\
	&\quad	-			\frac{1}{16}	\norm{ \sqrt{\tau} (U_t+ x U_y)}{T}^2
		+ \Big(	\beta -\frac{2\alpha^2}{3} - \frac{\delta_T \beta^2}{\alpha} -\delta_T \gamma\Big)\norm{U_y}{T}^2
\\	&\quad 	-	(\frac{3}{8} +2\delta_T \alpha) \norm{U_x}{T} ^2\bigg)
-\frac{1}{16}\norm{U}{\rm dG}^2,
	\end{aligned}
\]}
and, 	recalling $\alpha = (8\delta)^{-1}$, $\beta = (48\delta^2)^{-1}$ and $\gamma = 10(48^2\delta^3)^{-1}$, we calculate
	\begin{equation*}\begin{split}
		\beta -\frac{2\alpha^2}{3} - \frac{\delta_T \beta^2}{\alpha} -\delta_T \gamma
		=
		(384 \delta_T^2)^{-1} =&\ \frac{\alpha^2}{6}, \quad\text{and}\quad
		 \frac{3}{8 }  +2\delta_T \alpha =\ \frac{5}{8},
\end{split}\end{equation*}
on each $T\in\Th$.
The result follows upon using the last bound and the last two identities on \eqref{eq:inder_Vthree}.
\end{proof}

\end{document}
